\documentclass[11pt,aps,preprint,nofootinbib,superscriptaddress,tightenlines]{revtex4-2}
\usepackage{amsmath,amssymb,bm,mathtools,graphicx}
\usepackage{xcolor}
\definecolor{darkblue}{RGB}{0,0,160}
\usepackage[colorlinks=true,linkcolor=darkblue,citecolor=darkblue,urlcolor=darkblue]{hyperref}

\def\st{\begin{equation}}
\def\stp{\end{equation}}
\def\Eq#1{eq.~(\ref{#1})}
\def\Sect#1{Sect.~\ref{#1}}
\newcommand{\avg}[1]{\langle #1\rangle}
\let\oldparagraph\paragraph
\renewcommand{\paragraph}{\par\bigskip\oldparagraph}

\begin{document}

\title{The transverse rigidity of the particle arrangement: hypotheses and tests}
\author{D.~Teaney}
\affiliation{Department of Physics and Astronomy, Stony Brook University,
Stony Brook, New York 11794, USA}
\date{\today}

\begin{abstract}
The equilibrium runs (\texttt{campaigns.tex}) show transverse sound: the shear mode
oscillates with $c_t=0.022$ instead of diffusing.  This note separates what is measured
from what is conjectured, states the conjectures with the tests that can falsify them (all
in the campaign \texttt{rigidity}, now running), and lists ways to remove the rigidity, to be
discussed before any is implemented.
\end{abstract}

\maketitle

%=============================================================================
\section{What is measured}
\label{sec:facts}

All at $h/a=2$, Wendland C2 kernel, $\delta v=0.043$ ($s_a=541$), $\eta/s=0.2$
($\nu=\eta/w=2.2\times10^{-3}$).
\begin{enumerate}
\item $C_T(k,\tau)$ is fitted, within its errors, by the velocity correlator of a damped
  oscillator,
  \st
  C_T(k,\tau)=A\,e^{-\gamma\tau}\Big[\cos\omega\tau-\frac{\gamma}{\omega}\sin\omega\tau\Big],
  \qquad \omega^2+\gamma^2=c_t^2k^2,
  \label{osc}
  \stp
  with $c_t=0.022$ for $kh=0.13$--$0.74$, in boxes $1.2$ and $2.4$, from two starts and two
  versions of the code.  The fit recovers $c_t$ and $\gamma$ from synthetic data with known
  values, and returns $c_t=0$ for a synthetic fluid.
\item $2\gamma/\nu k^2=1.1$--$1.2$ at $kh\ge0.5$, rising to $\approx2$ at $kh=0.13$.
\item The local disorder (rms change of nearest-neighbour distances) saturates at $0.3a$: the
  particles do not exchange neighbours.
\item Ideal SPH alone (no viscous step), box $0.6$: from the lattice with thermal momenta
  the kinetic energy grows from $0.9$ to $6.8$ times its equilibrium value by $t=20$, still
  rising, while $H$ is conserved to $10^{-7}$.
\item Ideal SPH alone from random positions (smoothed entropy made uniform): the kinetic energy
  reaches $420$ times its equilibrium value by $t=5$.  The full scheme from random positions
  heats by $64\%$ in $H$ by $t=10$.
\end{enumerate}

%=============================================================================
\section{Hypothesis A: a broken relabeling symmetry gives a shear modulus}
\label{sec:A}

\paragraph{The continuum.}  The energy of a perfect fluid depends on the density and entropy
fields only.  A volume-preserving relabeling of the fluid elements leaves the fields, and so
the energy, unchanged: a transverse displacement $\bm u$ with $\nabla\cdot\bm u=0$ costs no
energy, and there is no shear modulus.

\paragraph{SPH.}  The smoothed entropy at particle $a$ is $S_a=\sum_bs_bW(\bm x_a-\bm x_b)$.  On a
lattice with equal $s_b=sV$,
\st
S_a=s\sum_{\bm G}\tilde W(\bm G)\,e^{i\bm G\cdot\bm x_a},
\stp
a sum over the reciprocal lattice.  The term $\bm G=0$ is the exact density; the others, set by
$\tilde W$ at $|\bm G|\ge b=4\pi/(\sqrt3a)$, are the lattice error.  A shear at fixed area changes
the $|\bm G|$, hence the sum, at second order in the strain $\epsilon$: the energy rises by
$\tfrac12G\epsilon^2$ per volume with, by this estimate,
\st
\frac{G}{p}\sim C\,(bh)^2\,\big|\tilde W''(bh)\big|\sim C\,m(m+1)\,\big|\tilde W(bh)\big|,
\label{Gest}
\stp
where $\tilde W\propto(bh)^{-m}$ asymptotically and $C$ is a number of order the count of
shortest $\bm G$ (six).  This is an order-of-magnitude argument, not a calculation; in
particular the forces in SPH depend on $\nabla W$ as well, and the stochastic step may add to
or remove from $G$.

\paragraph{Predictions.}  $\tilde W$ at the first reciprocal vector:
\begin{center}
\begin{tabular}{lcccc}
\hline
kernel & $\avg{r^2}/h^2$ & $\tilde W(bh)$, $h/a=2$ & $2.5$ & $3$ \\ \hline
Wendland C2 & $0.56$ & $1.2\times10^{-4}$ & $3.7\times10^{-5}$ & $1.5\times10^{-5}$ \\
Wendland C4 & $0.42$ & $3.5\times10^{-5}$ & $7.6\times10^{-6}$ & $2.2\times10^{-6}$ \\
Wendland C6 & $0.34$ & $1.6\times10^{-5}$ & $2.3\times10^{-6}$ & $5.0\times10^{-7}$ \\
\hline
\end{tabular}
\end{center}
If $G\propto\tilde W(bh)$, then $c_t=(G/w)^{1/2}\propto\tilde W(bh)^{1/2}$, and from $c_t=0.022$
for C2 at $h/a=2$:
\begin{center}
\begin{tabular}{lcc}
\hline
run & $c_t$ if $G\propto\tilde W$ & $c_t$ if $G\propto\tilde W^2$ \\ \hline
C2, $h/a=2.5$ & $0.012$ & $0.0068$ \\
C2, $h/a=3$ & $0.0078$ & $0.0028$ \\
C4, $h/a=2$ & $0.012$ & $0.0064$ \\
C4, $h/a=2.5$ & $0.0055$ & $0.0014$ \\
C6, $h/a=2$ & $0.0081$ & $0.0030$ \\
\hline
\end{tabular}
\end{center}
\paragraph{Falsified if:} $c_t$ does not fall with $h/a$ and with the kernel order; or C6 at
$h/a=2$ and C2 at $h/a=3$ (equal $\tilde W(bh)$, different $h/a$) give different $c_t$; or the
rigidity is absent when the positions are frozen \emph{and} present only with the stochastic
step (then it is not a property of $H$).

\paragraph{Second test of the same quantity.}  The runs record the displacements $\bm u_a$
from the initial sites.  For an elastic medium $\avg{|\dot u_T|^2}=c_t^2k^2\avg{|u_T|^2}$, and
$\bm M_{\bm k}\simeq wV\dot{\bm u}_{\bm k}$, so the statics alone give
\st
c_t^2=\frac{\avg{|M_T(\bm k)|^2}}{(wV)^2k^2\avg{|u_T(\bm k)|^2}} .
\stp
In a fluid $\avg{|u_T|^2}$ grows linearly in time without bound.  Agreement with the fit of
\Eq{osc} is a check that the oscillation is elastic.

%=============================================================================
\section{The ideal Hamiltonian's phonons (static test)}
\label{sec:phonons}

\texttt{examples/ideal\_lattice\_modes.jl} (branch \texttt{rigidity}).  On the perfect hexagonal
lattice (box $1.2$, fixed $s_a$, zero momenta), displace the particles from their sites by one
plane wave,
\st
\bm u_a\equiv\bm x_a-\bm x_a^0=\epsilon\,\hat{\bm e}\cos(\bm q\cdot\bm x_a^0),\qquad
\hat{\bm e}\parallel\bm q\ \text{(L)}\ \text{or}\ \perp\bm q\ \text{(T)},\qquad \epsilon=0.01a,
\stp
re-smooth, and take the second difference of $U=H(\{\bm x\})$:
\st
\Delta U=\tfrac12\big[U(\epsilon)+U(-\epsilon)-2U(0)\big]=\tfrac12\epsilon^2\,\frac N2\,K(\bm q),
\qquad \omega^2(\bm q)=\frac{K(\bm q)}{wV} .
\stp
This is the phonon spectrum of the ideal step, from the energy alone: no time stepping, no
noise.  A fluid has $\omega_L=c_sq$ and $\omega_T=0$.  Results ($c_t=\omega_T/q$ at
$qh=0.26$--$0.37$; it is constant there):
\begin{center}
\begin{tabular}{lccc}
\hline
kernel & $h/a=2$ & $2.5$ & $3$ \\ \hline
C2 & $c_t=0.034$ & $0.031$ & $0.022$ \\
C4 & $\omega_T^2<0$ (unstable) & unstable & unstable \\
C6 & $0.017$ & $0.0066$ & $0.0026$ \\
\hline
\end{tabular}
\end{center}
In every case $\omega_L/c_sq=1.000\pm0.003$ at $qh\le0.37$: the ideal step has the right sound
speed.
\begin{enumerate}
\item \emph{The rigidity is a property of $H$.}  C2 at $h/a=2$ has a linear transverse branch
  with $c_t=0.034$ on the perfect lattice; the runs give $0.022$, the same effect softened by
  the thermal disorder.
\item \emph{The simple estimate of \Sect{sec:A} fails quantitatively.}  For C2, $c_t$ falls by
  only $1.6$ from $h/a=2$ to $3$ while $\tilde W(bh)$ falls by $8$; for C4 the sign of the
  shear modulus is \emph{negative}: the hexagonal lattice is unstable to shear.  The shear
  modulus is a lattice sum whose terms change sign, not a single $\tilde W(bh)$.  C6 falls fast,
  a factor $6.5$ from $h/a=2$ to $3$, close to $\tilde W(bh)^{1/2}$.
\item \emph{Zero modes.}  For C2 at $h/a=2$ (all 648 shells up to $qh=7.5$), the only
  vanishing frequencies are at the zone boundary, $qh=7.255$ ($q=b/2$, neighbouring rows moving
  in opposite directions): $\omega_L=0$ and $\omega_T^2\approx-10^{-6}$.  These modes do not
  change $H$ to second order, so nothing restores them.  They are a candidate for the growth of
  the kinetic energy in ideal SPH (fact 4: energy flowing into neutral modes through the
  anharmonic terms) and for what the $\lambda=1$ noise drives into the melt.
\end{enumerate}
Consequences: (i) $G$ can be positive or negative depending on kernel and $h/a$, so it can be
made small, perhaps zero, on the perfect lattice by the choice of kernel and $h/a$; but the
relevant state is the thermally disordered one, where $c_t$ is smaller (0.022 against 0.034).
(ii) The C4 runs now going start on a lattice that is unstable to shear.  (iii) Any revision
of the ideal step, including $H=\sum_ae_a$, should be checked first with this static test:
$\omega_L=c_sq$, $\omega_T\to0$, and no zero or negative modes in the zone.

%=============================================================================
\section{Hypothesis B: the arrangement carries energy}
\label{sec:B}

Facts 4 and 5 show that at fixed $s_a$ the SPH energy depends on the arrangement, by a lot for
a random one ($\sim400T$ per particle released) and by a little for the lattice ($\sim6T$ per
particle by $t=20$, still growing).  The conjecture is that at temperature $T$ the equilibrium
of the positions lies near a minimum of this energy, so the particles form a solid whatever
the start, and a random arrangement is not an equilibrium state.

This is \emph{not} established.  Fact 4 says that the hexagonal lattice is not itself the
minimum (ideal SPH moves away from it, releasing energy), so the relevant minimum is some
nearby disordered arrangement, and why the full scheme stays at a disorder of $0.3a$ is not
explained.  The ideal-SPH runs at $h/a=2$, $2.5$, $3$ and for C4, C6 (to $t=50$) show whether
the kinetic energy saturates, and how the released energy scales with $\tilde W(bh)$.  If it
scales like $G$, A and B are the same effect.

%=============================================================================
\section{Ways out, if A holds}
\label{sec:ways}

\paragraph{Smaller lattice error.}  A fluid needs $c_tk\ll\nu k^2/2$ at the smallest $k$ of
interest,
\st
c_t\ll\frac{\nu k_{\rm min}}{2}=\frac{\pi\nu}{L},
\stp
about $3\times10^{-3}$ for $L=2.4$: a factor $7$ in $c_t$, $50$ in $G$.  Whether a kernel and
$h/a$ reach this at acceptable cost is what the runs measure.

\paragraph{Noise in $d\bm x/dt$.}  Particles that exchange neighbours relax shear stress in a
Maxwell time $\tau_M\sim a^2/D_x$ (diffusion constant $D_x$ of the positions relative to the
fluid), turning the solid into a fluid with an added viscosity $G\tau_M$.  This is small
compared with $\eta$ when
\st
D_x\gg\frac{c_t^2a^2}{\nu}\approx0.06\,\nu .
\stp
Three versions, with their expected outcomes:
\begin{enumerate}
\item \emph{Displacements alone} ($s_a$, $\bm m_a$ fixed): each move raises the arrangement
  energy (Hypothesis B), which ends up as heat; expected to heat strongly.
\item \emph{Displacements with a Metropolis test on $e^{S}$}: samples the true equilibrium of
  the positions; if B holds that is the solid, and $G$ is unchanged.
\item \emph{Relabeling}: displace a particle and refit $s_a$, $\bm m_a$, $e_a$ so that the
  smoothed fields are unchanged, as molecular self-diffusion leaves the hydrodynamic fields
  unchanged.  The energy change is then only the error of the refit, which a Metropolis test
  can absorb.  This restores the relabeling symmetry to the accuracy of the refit.  It needs a
  local linear solve per move that conserves $E$ and $\bm M$, and a check that it does not
  spoil the detailed balance of the stochastic step.
\end{enumerate}
A quick test of each at $L=0.6$ (heating and $c_t$ against $D_x$) would settle 1 and 2.

\paragraph{A different Hamiltonian: $H=\sum_ae_a$ with a partition of space (owner's suggestion).}
Write the energy as the sum of the particle energies,
\st
H=\sum_ae_a(s_a,\bm m_a,V_a),\qquad \sum_aV_a=A\ \text{exactly},
\label{Hpart}
\stp
with the volumes $V_a(\{\bm x\})$ a partition of the area $A$: Voronoi cells, or the partition of
unity $\psi_a(\bm x)=W(\bm x-\bm x_a)/\sum_bW(\bm x-\bm x_b)$, $V_a=\int\psi_a\,d^2x$, of the
meshless finite-mass methods.  Then:
\begin{enumerate}
\item \emph{No shear modulus for homogeneous strain (exact).}  An area-preserving affine map takes
  a Bravais lattice into another Bravais lattice, on which all cells are equivalent, so every
  $V_a$ stays $A/N$ and $H$ is unchanged: $G=0$.  The present kernel sum
  $S_a=\sum_bs_bW(\bm x_a-\bm x_b)$ has no such property; its change under shear is the lattice
  error of \Sect{sec:A}.
\item \emph{Conjecture:} non-affine motions change the $V_a$ at second order in the strain
  gradients, so the effective modulus vanishes as $(ka)^2$ at long wavelength, which is what a
  fluid needs.  Test: $c_t(k)$ in a run with \Eq{Hpart}; it should fall linearly with $k$.
\item \emph{One energy.}  The stochastic step conserves $\sum_ae_a$ exactly; with \Eq{Hpart} the
  Hamiltonian step conserves the same quantity, and the mismatch between the smoothed $H$ and
  $\sum_ae_a$, a candidate source of the slow heating, is gone by construction.
\end{enumerate}
Costs: the particles move with their own velocities, $\dot{\bm x}_a=\partial H/\partial\bm m_a$,
not smoothed ones, so the lattice-scale modes are less filtered; the forces need
$\partial V_b/\partial\bm x_a$ (Voronoi geometry or the partition-of-unity integrals).  And with
$e_a(s_a,V_a)$ each particle is its own thermodynamic system: the pressure and the equilibrium
distribution $e^{\sum_as_a}$ live at the scale $a$, and how the coarse-grained pressure and
entropy follow from them has to be worked out (the factor $9/16$ in the energy variance of
\texttt{eqstart.tex} is an instance).

\paragraph{Not proposed.}  Lowering $s_a$ to melt the arrangement: the scheme blows up at
$\delta v=0.14$ (\texttt{campaigns.tex}).

\end{document}
